\chapter{Job arrays: a parameter sweep}
\label{ch:arrays}

Say you want to run one program on hundreds of inputs. Submitting hundreds of separate jobs is clumsy. A \textbf{job array} submits them as one, with each \textbf{task} told its number (SchedMD, n.d.).

This chapter builds a working sweep. Download the three files: \href{examples/heat_loss.py}{\texttt{heat\_loss.py}}, \href{examples/make_cases.py}{\texttt{make\_cases.py}} and \href{examples/run_sweep.sh}{\texttt{run\_sweep.sh}}. The "simulation" estimates a building\textquotesingle s yearly heating need, but you can swap in your own program.

\section{The science in one paragraph}\label{the-science-in-one-paragraph}

Heat flows out through a wall in proportion to its area, how leaky it is (its \textbf{U-value}, in W per m² per kelvin) and the temperature difference. Air leaking out takes heat too: warming 1 m³ of air by 1 K takes about 0.33 Wh. \textbf{Degree days} measure how cold a year is, summing how far each day falls below a base temperature. So heat in kWh per year is:

Q = (U × A + 0.33 × n × V) × DD × 24 / 1000

where A is wall area, n is air changes per hour, V is volume and DD is degree days.

\section{Step 1: the program}\label{step-1-the-program}

\textbf{heat\_loss.py}

\lstinputlisting{../../../teaching/hpc/examples/heat_loss.py}

\section{Step 2: make the inputs}\label{step-2-make-the-inputs}

\textbf{make\_cases.py}

\lstinputlisting{../../../teaching/hpc/examples/make_cases.py}

Three settings, each with three values, give 3 × 3 × 3 = 27 combinations.

\begin{lstlisting}[language=bash]
# cluster, in a project folder holding the three files

# Find your Python first. Look for a module with: module avail 2>&1 | grep -i python
module load <python-module>
python3 make_cases.py cases
ls cases | head                # case_000.json, case_001.json ...
python3 heat_loss.py cases/case_000.json && cat cases/case_000.out    # 11016.0
\end{lstlisting}

\textbf{What you should see}

\begin{lstlisting}[language=text]
$ python3 make_cases.py cases
wrote 27 cases to cases

$ ls cases | head
case_000.json
case_001.json
case_002.json
case_003.json
case_004.json
case_005.json
case_006.json
case_007.json
case_008.json
case_009.json

$ python3 heat_loss.py cases/case_000.json && cat cases/case_000.out
11016.0
\end{lstlisting}

The last line is your test. For the first case, U = 2.1, A = 90, n = 1, V = 200 and DD = 1800, so by hand: (189 + 66) × 1800 × 24 / 1000 = 11016. Then delete that output, with \texttt{rm\ cases/case\_000.out}, so that the job below runs every case.

\section{Step 3: the job script}\label{step-3-the-job-script}

Running one input per task works, but very short tasks are wasteful, because starting each one has a cost. So \textbf{pack} several runs into each task.

\textbf{run\_sweep.sh}

\lstinputlisting{../../../teaching/hpc/examples/run_sweep.sh}

Add your \texttt{module\ load} line just before the \texttt{mapfile} line. Slurm usually passes your session\textquotesingle s environment on to the job, but stating it in the script is what makes the job repeatable.

\section{Step 4: how many tasks?}\label{step-4-how-many-tasks}

With N cases and P per task, divide and round up, so that the last few cases are covered. Tasks are numbered from 0. For 27 cases at 5 per task, that is 6 tasks, numbered 0 to 5.

\begin{lstlisting}[language=bash]
# cluster

# Line endings, if you edited on Windows
sed -i 's/\r$//' run_sweep.sh

# The log folder must exist before submission
mkdir -p logs

# Count the cases, and work out the last task number
N=$(ls cases/case_*.json | wc -l)
PER=5
# In whole-number arithmetic, (N + P - 1) / P rounds up
LAST=$(( (N + PER - 1) / PER - 1 ))
echo "$N cases, tasks 0 to $LAST"

sbatch --array=0-$LAST run_sweep.sh cases $PER
\end{lstlisting}

\textbf{What you should see from the echo line}

\begin{lstlisting}[language=text]
27 cases, tasks 0 to 5
\end{lstlisting}

The \texttt{sbatch} line then prints the number of your job, which differs each time.

A useful extra: \texttt{-\/-array=0-99\%10} runs at most 10 tasks at once, which is polite on a busy system.

\section{Step 5: did it all work?}\label{step-5-did-it-all-work}

\begin{lstlisting}[language=bash]
# cluster
squeue -u $USER                       # wait until it is empty
ls cases/*.out | wc -l                # expect 27
grep -l FAILED logs/*.out | wc -l     # expect 0
tail -2 logs/sweep_*_0.out            # read one log
\end{lstlisting}

If some failed, resubmit the same command. Because each task skips cases that already have an output, only the missing ones run. Being safe to re-run like this makes large runs much less stressful.

\section{Step 6: collect the results}\label{step-6-collect-the-results}

\begin{lstlisting}[language=bash]
# cluster
# A loop that prints each case's name and answer into a CSV file.
# ${f%.out} strips the ending, and $(cat "$f") inserts the file's contents
for f in cases/*.out; do echo "$(basename "${f%.out}"),$(cat "$f")"; done > results.csv
head -3 results.csv
\end{lstlisting}

\textbf{What you should see from head}

\begin{lstlisting}[language=text]
case_000,11016.0
case_001,12240.0
case_002,13464.0
\end{lstlisting}

\begin{lstlisting}[language=powershell]
# laptop
scp kelvin:~/project/results.csv "D:\My Work"
\end{lstlisting}

\section{Choosing a packing}\label{choosing-a-packing}

{\def\LTcaptype{none} % do not increment counter
\begin{longtable}[]{@{}
  >{\raggedright\arraybackslash}p{(\linewidth - 4\tabcolsep) * \real{0.3333}}
  >{\raggedright\arraybackslash}p{(\linewidth - 4\tabcolsep) * \real{0.3333}}
  >{\raggedright\arraybackslash}p{(\linewidth - 4\tabcolsep) * \real{0.3333}}@{}}
\toprule\noalign{}
\begin{minipage}[b]{\linewidth}\raggedright
\end{minipage} & \begin{minipage}[b]{\linewidth}\raggedright
Few runs per task (many tasks)
\end{minipage} & \begin{minipage}[b]{\linewidth}\raggedright
Many runs per task (few tasks)
\end{minipage} \\
\midrule\noalign{}
\endhead
\bottomrule\noalign{}
\endlastfoot
Finishes & Sooner, if cores are free & Later: each task takes longer \\
Overhead & Higher: more starts, more load on the scheduler & Lower \\
After a failure & Less lost & More to redo, unless it is restart-safe \\
\end{longtable}
}

There is no single right answer. Test with your own program, with care: submit each setting separately, repeat it, and keep the copying of files out of your timing. If you are planning a very large run, ask your support team what they prefer.

